% Results Section for Non-Convex Polygon Packing Study
\section{Experimental Results}

\subsection{Experimental Setup and Hardware Constraints}
All experiments were conducted under a strict 20 CPU core limit. The solver was compiled via standard GCC toolchains. We evaluated three parallel coordination strategies across a graph suite of $N \in \{5, 10, 20, 50, 100\}$ using 10 threads per run and 2 concurrent seeds, ensuring exactly 20 CPUs were utilized.

\subsection{Calibration and Gate Verification}
Prior to the full suite, all methods passed preliminary gating for deadlock prevention, bracket shrinking, and baseline feasibility.

\begin{itemize}
  \item \textbf{PT Pilot Calibration:} The Parallel Tempering temperature ladder was tuned at $N=100$. The final swap acceptance rate was recorded at \textbf{[XX.X]\%}, falling within our target operational bounds of $[15\%, 60\%]$.
  \item \textbf{Gate Status:} All methods successfully shrunk the bisection bracket by $\geq 3$ probes at $N=20$ (Gate A) and achieved baseline feasibility within a 10-minute budget at $N=100$ (Gate C).
\end{itemize}

\subsection{Algorithm Comparison (The Graph Suite)}
The graph suite evaluated Multi-Start (MS), Elite Resample Multi-Start (ER-MS), and Parallel Tempering (PT).

\subsubsection{Convergence and Packing Efficiency}
The primary metric for comparison is the best feasible bounding box side length, $L_{\text{best}}(t)$, and the resulting packing efficiency, $\eta$.

\begin{table}[h]
\centering
\begin{tabular}{c|c|c|c|c}
$N$ & MS Best $L$ & ER-MS Best $L$ & PT Best $L$ & Best Overall Efficiency ($\eta$) \\
\hline
5   & [Value]     & [Value]        & [Value]     & [Value]\% \\
10  & [Value]     & [Value]        & [Value]     & [Value]\% \\
20  & [Value]     & [Value]        & [Value]     & [Value]\% \\
50  & [Value]     & [Value]        & [Value]     & [Value]\% \\
100 & [Value]     & [Value]        & [Value]     & [Value]\% \\
\end{tabular}
\caption{Best $L$ and packing efficiency $\eta$ for each method and $N$.}
\end{table}

% Insert L_vs_N.png and eta_vs_N.png here

\subsubsection{Method Selection}
Based on the convergence plots, \textbf{[Insert Winning Method]} demonstrated the most consistent bracket reduction and highest packing density across $N \geq 50$. Specifically, the mechanism of \textbf{[e.g., elite resampling / geometric replica exchange]} allowed it to avoid the premature convergence observed in \textbf{[Losing Method]}. Therefore, \textbf{[Winning Method]} was selected for the $N=200$ Hero Run.

\subsection{Hero Run Demonstration ($N=200$)}
A singular demonstration run was executed at $N=200$ using \textbf{[Winning Method]} utilizing all 20 available CPU cores. The run consisted of a 150-minute bisection phase followed by a 50-minute polish phase at a fixed $L_{\text{best}}$.

\subsubsection{Hero Run Performance Metrics}
\begin{itemize}
  \item \textbf{Total Wall-Clock Time:} [XX] hours [XX] minutes
  \item \textbf{Final Bounding Box ($L_{\text{best}}$):} [Value]
  \item \textbf{Final Bracket Width ($L_{hi} - L_{lo}$):} [Value]
  \item \textbf{Packing Efficiency ($\eta$):} [Value]\%
\end{itemize}

% Insert N200_best_state.svg here

\subsubsection{Polish Phase Behavior}
During the 50-minute polish phase, the thermal pulse rule (temperature $\times 10$ for 30 seconds after 10 minutes of stagnation) was triggered \textbf{[X]} times. This successfully perturbed the system, yielding a final $L$ reduction of \textbf{[Value]} compared to the end of the strict bisection phase.
